\capitulo{5}{Aspectos relevantes del desarrollo del proyecto}

En este capítulo se describen los aspectos más relevantes y representativos del
desarrollo del proyecto, poniendo el foco en las decisiones técnicas y de
diseño adoptadas a lo largo de su ciclo de vida.
\section{Arquitectura y diseño del software}

Durante la fase de análisis y diseño del proyecto se tomaron decisiones clave.
Debido a la naturaleza de los datos de entrada, documentos y textos que se
obtienen de Github, provenientes de \textit{issues}, \textit{readmes}, y
documentos en LaTeX de los distintos \acrshort{tfg}, se decidió implementar un
proceso específico de limpieza y normalización del texto, desacoplado del resto
del sistema para facilitar su reutilización y evolución. También debido a la
necesidad de experimentar con distintos enfoques de modelado de temas, se optó
por diseñar una interfaz que permitiera la abstracción y asi integrar
diferentes proveedores de modelos. Estas necesidades influyeron directamente en
el diseño del sistema, requiriendo una arquitectura flexible que permitiera
integrar múltiples modelos sin modificar la lógica principal de la aplicación.

Durante la fase de análisis se identificaron claramente tres responsabilidades
principales, que dieron lugar a la división del sistema en módulos
independientes:

\begin{itemize}
      \item Un módulo de \textbf{ingestión}, responsable de la obtención y normalización
            inicial de los documentos fuente.
      \item Un módulo de \textbf{procesamiento}, encargado del preprocesado del texto y del
            modelado de tópicos.
      \item Un módulo de \textbf{visualización}, orientado a la exploración y análisis de
            los resultados generados.
\end{itemize}

Para el diseño del código interno de cada modulo se empleó una arquitectura de
tipo hexagonal~\cite{HexagonalArchitectureClean2022}. Este enfoque permite
desacoplar la lógica de negocio de los detalles de infraestructura, como el
acceso a datos, la ejecución de modelos o la interfaz de usuario
(\acrfull{ui}). Esta decisión permitió separar claramente la lógica de dominio
del procesamiento de tópicos de los detalles de infraestructura, como el acceso
a datos, la ejecución de modelos concretos o la interfaz de usuario.

Además, la arquitectura adoptada mejoró la testabilidad del sistema y
contribuyó a una organización del código más clara y mantenible, evitando un
crecimiento desordenado del proyecto a medida que se incorporaban nuevas
funcionalidades.

\section{Gestión de los datos y almacenamiento de resultados}

Uno de los aspectos relevantes del desarrollo del proyecto fue la elección del
formato de almacenamiento de los datos intermedios y de los resultados
generados por los distintos modelos. Para este propósito se optó por utilizar
el formato \textbf{Parquet}~\cite{ParquetDocumentation}, un formato de
almacenamiento en columnas ampliamente utilizado en entornos de análisis de
datos y \textit{big data}.

La elección de Parquet frente a alternativas más simples como \acrfull{csv} o
formatos binarios genéricos se basó en varios criterios. En primer lugar,
Parquet ofrece un almacenamiento más eficiente en espacio como en tiempo de
acceso, gracias a su organización en columnas y a los mecanismos de compresión
integrados. Este aspecto resulta especialmente relevante en un proyecto que
genera múltiples ejecuciones de modelos, cada una con diferentes
configuraciones de hiperparámetros y métricas asociadas.

Desde el punto de vista de la implementación, el formato Parquet se utilizó
principalmente para almacenar:
\begin{itemize}
      \item Los datos sin procesar, incluyendo los textos y cualquier metadato relevante.
      \item Los resultados de la ejecución de los modelos de temas, incluyendo la
            información asociada a los temas generados.
      \item Las métricas de evaluación obtenidas en cada ejecución.
      \item Los valores de los hiperparámetros utilizados durante los procesos de
            optimización.
\end{itemize}

Estos ficheros se organizan siguiendo una estructura de directorios coherente
con los módulos del sistema, lo que facilita la localización y reutilización de
los resultados generados. Además, el uso de un formato estándar y ampliamente
soportado permite que los datos puedan ser consumidos directamente tanto por el
módulo de visualización como por herramientas externas de análisis, como
entornos interactivos de experimentación.

\section{Optimización de hiperparámetros y gestión de experimentos}

Uno de los aspectos más relevantes de la implementación fue la integración de
Optuna para la optimización de hiperparámetros. Frente a herramientas más
genéricas como \textit{GridSearchCV} de \textit{scikit-learn}, Optuna no impone
restricciones sobre la estructura del modelo ni sobre el flujo de
entrenamiento, lo que facilita su integración con bibliotecas heterogéneas como
LDA, BERTopic, FASTopic o Top2Vec. Esta independencia del \textit{framework}
resulta clave en un proyecto que combina modelos con paradigmas y requisitos
muy diferentes.

Los resultados de cada ejecución, junto con los parámetros utilizados y las
métricas calculadas, se almacenan en formato Parquet. Esta decisión facilita la
trazabilidad de los experimentos, permite analizar ejecuciones pasadas y evita
la repetición innecesaria de cálculos costosos desde el punto de vista
computacional.

\section{Interfaz web y exploración de resultados}

El módulo de visualización se implementó como un frontal web utilizando
Streamlit. Este componente consume directamente los resultados generados por el
módulo de procesamiento y almacenados en ficheros Parquet, manteniendo así un
acoplamiento mínimo entre ambos módulos.

La interfaz permite seleccionar modelos, explorar tópicos generados, visualizar
métricas y analizar resultados de forma interactiva. Esta decisión sustituye
salidas estáticas por una herramienta más adecuada para un análisis
exploratorio y facilita tanto la evaluación del sistema como la comprensión de
los resultados obtenidos.